% !TeX program = lualatex
% ================================================================
% Polish Parallel Text Version of DevelopingAveragesStarters
%
% Generated by the server using the following settings:
%
%   language            Polish (pol)
%   text direction      left to right
%   font                Noto Serif
%   fallback font       Noto Serif
%   built from          latex/starters/targeted/KS3/statistics/developingAveragesStarters.tex
%   vocabulary key      latex/starters/targeted/KS3/statistics/developingAveragesStarters/L2/pol/developingAveragesStarters_polishKey.tex
%   tier 3 vocabulary   green   RGB 0 128 0
%   English text        black   RGB 0 0 0
%   Polish text         purple  RGB 112 48 160
%   layout macros       version 3
%
% Please don't edit any information in this comment block - the server
% compares it against the current settings and rebuilds this file when
% they differ, so changes here would be overwritten. However, feel free
% to make updates to the LaTeX below if you want to edit it.
% ================================================================

\documentclass{beamer}
% ================================================================
% Copied from the English sheet this was translated from, so that
% anything it sets up for itself still works here
% ================================================================
\usepackage{tikz}
\usetikzlibrary{calc,arrows.meta,patterns}
\usepackage{multicol}
\usepackage{amsmath}

% ================================================================
% Answer-overlay helpers
% ================================================================
\newcommand{\ablank}[1]{%
  \alt<2>{\textcolor{red}{#1}}{\underline{\phantom{#1}}}%
}
\newcommand{\ashow}[1]{\uncover<2->{\textcolor{red}{\small #1}}}
\newcommand{\ashowq}[1]{\alt<2->{\textcolor{red}{\small #1}}{?}}

% Answers drawn straight on to a TikZ diagram, revealed on overlay 2 like \ashow.
% Space is reserved on overlay 1, so the diagram does not move between slides.
\usetikzlibrary{overlay-beamer-styles}
\tikzset{
  ans/.style     = {red, thick, visible on=<2->},
  ansfill/.style = {red!15, visible on=<2->},
  anslab/.style  = {red, font=\tiny, visible on=<2->},
}

\title{Averages: Mean \& Median}
\author{}
\date{}

% Lego tower: \legotower{x-offset}{height}{colour}
% Each brick is 0.65 wide x 0.55 tall, with two studs
\newcommand{\legotower}[3]{%
  \foreach \i in {1,...,#2}{%
    \draw[fill=#3!55, draw=#3!80!black, line width=0.35pt]
      (#1, {(\i-1)*0.55}) rectangle (#1+0.65, {\i*0.55});
    \draw[fill=#3!80, draw=#3!80!black, line width=0.2pt]
      (#1+0.15,{(\i-1)*0.55+0.38}) circle (0.08);
    \draw[fill=#3!80, draw=#3!80!black, line width=0.2pt]
      (#1+0.44,{(\i-1)*0.55+0.38}) circle (0.08);
  }%
}

% Characters this language's font has not got are borrowed from another,
% rather than being silently dropped - see docs/translated-sheet-latex.md
\directlua{%
  if luaotfload and luaotfload.add_fallback then
    luaotfload.add_fallback("eallatinfallback",
      { "Noto Serif:mode=harf;" })
  else
    tex.error("this luaotfload is too old to register a fallback font")
  end
}
\usepackage[english]{babel}
\babelprovide[import]{polish}
\babelfont[polish]{rm}[Renderer=Harfbuzz, BoldFont={Noto Serif}, ItalicFont={Noto Serif}, BoldItalicFont={Noto Serif}, RawFeature={fallback=eallatinfallback}]{Noto Serif}
\babelfont[polish]{sf}[Renderer=Harfbuzz, BoldFont={Noto Serif}, ItalicFont={Noto Serif}, BoldItalicFont={Noto Serif}, RawFeature={fallback=eallatinfallback}]{Noto Serif}
\newcommand{\ealtext}[1]{\foreignlanguage{polish}{#1}}
\newcommand{\ealtextblock}[1]{{\raggedright\foreignlanguage{polish}{#1}\par}}
% ================================================================
% EAL layout helpers
% ================================================================
\usepackage{xcolor}

\definecolor{ealkeycolour}{RGB}{0,128,0}
\definecolor{ealenglish}{RGB}{0,0,0}
\definecolor{ealtwocolour}{RGB}{112,48,160}

\newsavebox{\ealboxen}
\newsavebox{\ealboxtr}
\newlength{\ealglosswd}
\newlength{\ealglossraise}
\setlength{\ealglossraise}{1.15em}

% a tier 3 word, in English and in translation alike
\newcommand{\ealkey}[1]{\textcolor{ealkeycolour}{#1}}
\newcommand{\ealkeytr}[1]{\textcolor{ealkeycolour}{\ealtext{#1}}}

% One translated word above one English word. Built from kernel
% primitives, never a tabular, which would break wherever the sheet
% loads array, booktabs or tabularx. The box is as wide as the wider of
% the two, so a long translation pushes its neighbours aside instead of
% overlapping them, and it claims its full height so a table row or a
% TikZ node makes room for it.
\newcommand{\ealgloss}[2]{%
  \sbox{\ealboxen}{\ealkey{#1}}%
  \sbox{\ealboxtr}{\small\textcolor{ealtwocolour}{\ealtext{#2}}}%
  \setlength{\ealglosswd}{\wd\ealboxen}%
  \ifdim\wd\ealboxtr>\ealglosswd\setlength{\ealglosswd}{\wd\ealboxtr}\fi
  \leavevmode
  \makebox[\ealglosswd][c]{%
    \makebox[0pt][c]{\usebox{\ealboxen}}%
    \makebox[0pt][c]{\raisebox{\ealglossraise}{\usebox{\ealboxtr}}}%
  }%
}

% opens up line spacing around a block of glosses, so the raised
% translations sit clear of the line above
\newenvironment{ealglossed}{\par\linespread{2.1}\selectfont\ignorespaces}{\par}

% a whole translated sentence above its English counterpart. block level,
% so it does not belong inside a tabular cell
\newcommand{\ealpara}[2]{%
  \par\addvspace{0.5\baselineskip}%
  {\color{ealtwocolour}\ealtextblock{#1}}%
  \penalty200
  {\color{ealenglish}#2\par}%
  \addvspace{0.5\baselineskip}%
}

\begin{document}

\begin{frame}{Starter}
\vspace{-0.8em}

\footnotesize
\begin{enumerate}\itemsep2pt
\begin{multicols}{2}
  \item $6 + 9 = \ashowq{15}$
  \item $24 \div 4 = \ashowq{6}$
  \item $7 + 13 + 5 = \ashowq{25}$
  \item $30 \div 5 = \ashowq{6}$
\end{multicols}
  \item \ealpara{Jak obliczamy \ealkeytr{średnią}?}{How do we work out the \ealkey{mean}?} 
        \ealpara{Dodaj \ealkeytr{wartości} i podziel przez ich liczbę.}{Add the \ealkey{values} and divide by how many there are.} 
        \ashow{Add the \ealkey{values} and divide by how many there are.}
  \item \ealpara{Jak obliczamy \ealkeytr{medianę}?}{How do we work out the \ealkey{median}?} 
        \ealpara{Uporządkuj \ealkeytr{dane} i znajdź środkową \ealkeytr{wartość}.}{Order the \ealkey{data} and find the middle \ealkey{value}.} 
        \ashow{Order the \ealkey{data} and find the middle \ealkey{value}.}
  \item \ealpara{Znajdź \ealkeytr{średnią} z:}{Find the \ealkey{mean} of:} $3,\ 7,\ 5,\ 9$ \ashow{$=6$}
  \item \ealpara{Znajdź \ealkeytr{medianę} z:}{Find the \ealkey{median} of:} $4,\ 9,\ 2,\ 7,\ 8$ \ashow{$=7$}
  \item \ealpara{Znajdź \ealkeytr{średnią} z:}{Find the \ealkey{mean} of:} $\frac{1}{2},\ \frac{3}{2},\ 2,\ 3$ \ashow{$=\frac{7}{4}$}
  \item \ealpara{Znajdź \ealkeytr{medianę} z:}{Find the \ealkey{median} of:} $-4,\ 2,\ -1,\ -7,\ 1$ \ashow{$=-1$}
  \item \ealpara{\ealkeytr{Średnia} z $x,\ x+4,\ x+8,\ x+12$ wynosi $15$, jaka jest suma wszystkich \ealkeytr{wartości}?}{The \ealkey{mean} of $x,\ x+4,\ x+8,\ x+12$ is $15$, what is the total of all the \ealkey{values}?} \ashow{$60$}
  \item \ealpara{Znajdź \ealkeytr{wartość} $x$ w poprzednim pytaniu}{Find the \ealkey{value} of $x$ in the previous question} \ashow{$x=9$}
\end{enumerate}

\end{frame}


\begin{frame}{Starter -- Solutions}
\vspace{-0.8em}

\footnotesize
\begin{enumerate}\itemsep2pt
\begin{multicols}{2}
  \item $6 + 9 = \mathord{?}$ \textcolor{red}{$15$}
  \item $24 \div 4 = \mathord{?}$ \textcolor{red}{$6$}
  \item $7 + 13 + 5 = \mathord{?}$ \textcolor{red}{$25$}
  \item $30 \div 5 = \mathord{?}$ \textcolor{red}{$6$}
\end{multicols}
  \item \ealpara{Jak obliczamy \ealkeytr{średnią}?}{How do we work out the \ealkey{mean}?} 
        \ealpara{Dodaj \ealkeytr{wartości} i podziel przez ich liczbę.}{Add the \ealkey{values} and divide by how many there are.} 
        \textcolor{red}{Add the \ealkey{values} and divide by how many there are}
  \item \ealpara{Jak obliczamy \ealkeytr{medianę}?}{How do we work out the \ealkey{median}?} 
        \ealpara{Uporządkuj \ealkeytr{dane} i znajdź środkową \ealkeytr{wartość}.}{Order the \ealkey{data} and find the middle \ealkey{value}.} 
        \textcolor{red}{Order the \ealkey{data} and find the middle \ealkey{value}}
  \item \ealpara{Znajdź \ealkeytr{średnią} z:}{Find the \ealkey{mean} of:} $3,\ 7,\ 5,\ 9$ \textcolor{red}{$=6$}
  \item \ealpara{Znajdź \ealkeytr{medianę} z:}{Find the \ealkey{median} of:} $4,\ 9,\ 2,\ 7,\ 8$ \textcolor{red}{$=7$}
  \item \ealpara{Znajdź \ealkeytr{średnią} z:}{Find the \ealkey{mean} of:} $\frac{1}{2},\ \frac{3}{2},\ 2,\ 3$ \textcolor{red}{$=\frac{7}{4}$}
  \item \ealpara{Znajdź \ealkeytr{medianę} z:}{Find the \ealkey{median} of:} $-4,\ 2,\ -1,\ -7,\ 1$ \textcolor{red}{$=-1$}
  \item \ealpara{\ealkeytr{Średnia} z $x,\ x+4,\ x+8,\ x+12$ wynosi $15$, jaka jest suma?}{The \ealkey{mean} of $x,\ x+4,\ x+8,\ x+12$ is $15$, what is the total?} \textcolor{red}{$60$}
  \item \ealpara{Znajdź \ealkeytr{wartość} $x$}{Find the \ealkey{value} of $x$} \textcolor{red}{$x=9$}
\end{enumerate}
\end{frame}


\begin{frame}{Starter}
\vspace{-0.8em}

\footnotesize
\begin{enumerate}\itemsep2pt
\begin{multicols}{2}
  \item $3 + 7 + 10 = \ashowq{20}$
  \item $48 \div 6 = \ashowq{8}$
  \item $\frac{1}{2} + \frac{1}{4} = \ashowq{\frac{3}{4}}$
  \item $28 \div 4 = \ashowq{7}$
\end{multicols}
  \item \ealpara{Która \ealkeytr{średnia} wymaga najpierw uporządkowania \ealkeytr{danych}?}{Which \ealkey{average} requires you to order the \ealkey{data} first?} 
        \ealpara{\ealkeytr{Mediana}}{\ealkey{Median}} 
        \ashow{\ealkey{Median}}
  \item \ealpara{Znajdź \ealkeytr{medianę} z:}{Find the \ealkey{median} of:} $7,\ 2,\ 9,\ 4,\ 5$ \ashow{$=5$}
  \item \ealpara{Znajdź \ealkeytr{medianę} z:}{Find the \ealkey{median} of:} $3,\ 8,\ 1,\ 6,\ 4,\ 10$ \ashow{$=5$}
  \item \ealpara{Znajdź \ealkeytr{średnią} i \ealkeytr{medianę} z:}{Find the \ealkey{mean} and \ealkey{median} of:} $2,\ 5,\ 3,\ 8,\ 7$. 
        \ealpara{Która jest większa?}{Which is larger?} 
        \ealpara{\ealkeytr{Średnia} $=5$, \ealkeytr{mediana} $=5$ (równa).}{\ealkey{Mean} $=5$, \ealkey{median} $=5$ (equal).}
        \ashow{\ealkey{Mean} $=5$, \ealkey{median} $=5$ (equal).}
  \item \ealpara{Znajdź \ealkeytr{średnią} i \ealkeytr{medianę} z:}{Find the \ealkey{mean} and \ealkey{median} of:} $-5,\ 1,\ -2,\ 4,\ 7$ 
        \ealpara{\ealkeytr{Średnia} $=1$, \ealkeytr{mediana} $=1$.}{\ealkey{Mean} $=1$, \ealkey{median} $=1$.}
        \ashow{\ealkey{Mean} $=1$, \ealkey{median} $=1$.}
  \item \ealpara{\ealkeytr{Średnia} z $a,\ a+3,\ a+6$ wynosi $6$.}{The \ealkey{mean} of $a,\ a+3,\ a+6$ is $6$.}
        \begin{itemize}\footnotesize
        \item \ealpara{Jaka jest \ealkeytr{mediana} wyrażona przez $a$?}{What is the \ealkey{median} in terms of $a$?} \ashow{$a+3$}
          \item \ealpara{Jaka jest suma \ealkeytr{danych} (użyj \ealkeytr{średniej})?}{What is the total of the \ealkey{data} (use the \ealkey{mean})} \ashow{$18$}
          \item \ealpara{Znajdź $a$}{Find $a$} \ashow{$a=3$}
        \end{itemize}
\end{enumerate}

\end{frame}


\begin{frame}{Starter -- Solutions}
\vspace{-0.8em}

\footnotesize
\begin{enumerate}\itemsep2pt
\begin{multicols}{2}
  \item $3 + 7 + 10 = \mathord{?}$ \textcolor{red}{$20$}
  \item $48 \div 6 = \mathord{?}$ \textcolor{red}{$8$}
  \item $\frac{1}{2} + \frac{1}{4} = \mathord{?}$ \textcolor{red}{$\frac{3}{4}$}
  \item $28 \div 4 = \mathord{?}$ \textcolor{red}{$7$}
\end{multicols}
  \item \ealpara{Która \ealkeytr{średnia} wymaga najpierw uporządkowania?}{Which \ealkey{average} requires ordering first?} 
        \textcolor{red}{\ealkey{Median}}
  \item \ealpara{\ealkeytr{Mediana} z $7,\ 2,\ 9,\ 4,\ 5$}{\ealkey{Median} of $7,\ 2,\ 9,\ 4,\ 5$} 
        \textcolor{red}{$=5$}
  \item \ealpara{\ealkeytr{Mediana} z $3,\ 8,\ 1,\ 6,\ 4,\ 10$}{\ealkey{Median} of $3,\ 8,\ 1,\ 6,\ 4,\ 10$} 
        \textcolor{red}{$=5$}
  \item \ealpara{\ealkeytr{Średnia} i \ealkeytr{mediana} z $2,\ 5,\ 3,\ 8,\ 7$}{\ealkey{Mean} and \ealkey{median} of $2,\ 5,\ 3,\ 8,\ 7$} 
        \textcolor{red}{\ealkey{Mean} $=5,\ \text{\ealkey{median} }=5$}
  \item \ealpara{\ealkeytr{Średnia} i \ealkeytr{mediana} z $-5,\ 1,\ -2,\ 4,\ 7$}{\ealkey{Mean} and \ealkey{median} of $-5,\ 1,\ -2,\ 4,\ 7$} 
        \textcolor{red}{\ealkey{Mean} $=1,\ \text{\ealkey{median} }=1$}
  \item \ealpara{\ealkeytr{Średnia} z $a,\ a+3,\ a+6$ wynosi $6$}{\ealkey{Mean} of $a,\ a+3,\ a+6$ is $6$}
        \begin{itemize}\footnotesize
          \item \ealpara{\ealkeytr{Mediana}}{\ealkey{Median}} \textcolor{red}{$=a+3$}
          \item \ealpara{Suma}{Total} \textcolor{red}{$=18$}
          \item $a=$ \textcolor{red}{$3$}
        \end{itemize}
\end{enumerate}
\end{frame}


\begin{frame}{Starter}
\vspace{-0.8em}

\footnotesize
\begin{enumerate}\itemsep2pt
\begin{multicols}{2}
  \item $4 + 11 = \ashowq{15}$
  \item $40 \div 8 = \ashowq{5}$
  \item $\frac{3}{4} - \frac{1}{4} = \ashowq{\frac{1}{2}}$
  \item $72 \div 9 = \ashowq{8}$
  \end{multicols}
  \item \ealpara{Właścicielka sklepu z butami mówi: „Najpopularniejszy rozmiar w zeszłym tygodniu to rozmiar 6”. Której \ealkeytr{średniej} używa?}{A shoe shop owner says, ``The most popular size last week was size 6.'' Which \ealkey{average} is she using?} 
        \ealpara{\ealkeytr{Dominanta}}{\ealkey{Mode}} 
        \ashow{\ealkey{Mode}}
  \item \ealpara{Znajdź \ealkeytr{dominantę} z:}{Find the \ealkey{mode} of:} $3,\ 7,\ 3,\ 9,\ 5,\ 3,\ 7$ \ashow{$3$}
  \item \ealpara{Znajdź \ealkeytr{dominantę} z:}{Find the \ealkey{mode} of:} $2,\ 5,\ 8,\ 5,\ 3,\ 8,\ 1$ \ashow{$5$ and $8$}
  \item \ealpara{Znajdź \ealkeytr{średnią} i \ealkeytr{medianę} z:}{Find the \ealkey{mean} and \ealkey{median} of:} $4,\ 7,\ 2,\ 9,\ 3$ 
        \ealpara{\ealkeytr{Średnia} $=5$, \ealkeytr{mediana} $=4$.}{\ealkey{Mean} $=5$, \ealkey{median} $=4$.}
        \ashow{\ealkey{Mean} $=5$, \ealkey{median} $=4$.}
  \item \ealpara{Znajdź \ealkeytr{średnią}, \ealkeytr{medianę} i \ealkeytr{dominantę} z:}{Find the \ealkey{mean}, \ealkey{median} and \ealkey{mode} of:} $5,\ 3,\ 7,\ 3,\ 5,\ 5,\ 9$ 
        \ealpara{\ealkeytr{Średnia} $=\frac{37}{7}$, \ealkeytr{mediana} $=5$, \ealkeytr{dominanta} $=5$.}{\ealkey{Mean} $=\frac{37}{7}$, \ealkey{median} $=5$, \ealkey{mode} $=5$.}
        \ashow{\ealkey{Mean} $=\frac{37}{7}$, \ealkey{median} $=5$, \ealkey{mode} $=5$.}
  \item \ealpara{\ealkeytr{Zbiór} \ealkeytr{danych} to $\{4,\ 7,\ 4,\ x,\ 7,\ 3\}$}{The \ealkey{data} \ealkey{set} is $\{4,\ 7,\ 4,\ x,\ 7,\ 3\}$}
        \ealpara{\textbf{Jedyną} \ealkeytr{dominantą} jest $4$.}{The \textbf{only} \ealkey{mode} is $4$.}
        \ealpara{Znajdź $x$ i wyjaśnij swoje rozumowanie.}{Find $x$ and explain your reasoning.}
        \ealpara{$x=4$; wtedy $4$ występuje trzy razy.}{$x=4$; then $4$ occurs three times.}
        \ashow{$x=4$; then $4$ occurs three times.}
\end{enumerate}

\end{frame}


\begin{frame}{Starter -- Solutions}
\vspace{-0.8em}

\footnotesize
\begin{enumerate}\itemsep2pt
\begin{multicols}{2}
  \item $4 + 11 = \mathord{?}$ \textcolor{red}{$15$}
  \item $40 \div 8 = \mathord{?}$ \textcolor{red}{$5$}
  \item $\frac{3}{4} - \frac{1}{4} = \mathord{?}$ \textcolor{red}{$\frac{1}{2}$}
  \item $72 \div 9 = \mathord{?}$ \textcolor{red}{$8$}
\end{multicols}
  \item \ealpara{Która \ealkeytr{średnia} jest używana?}{Which \ealkey{average} is used?} 
        \textcolor{red}{\ealkey{Mode}}
  \item \ealpara{\ealkeytr{Dominanta} z $3,\ 7,\ 3,\ 9,\ 5,\ 3,\ 7$}{\ealkey{Mode} of $3,\ 7,\ 3,\ 9,\ 5,\ 3,\ 7$} 
        \textcolor{red}{$3$}
  \item \ealpara{\ealkeytr{Dominanta} z $2,\ 5,\ 8,\ 5,\ 3,\ 8,\ 1$}{\ealkey{Mode} of $2,\ 5,\ 8,\ 5,\ 3,\ 8,\ 1$} 
        \textcolor{red}{$5\ \text{and}\ 8$}
  \item \ealpara{\ealkeytr{Średnia} i \ealkeytr{mediana} z $4,\ 7,\ 2,\ 9,\ 3$}{\ealkey{Mean} and \ealkey{median} of $4,\ 7,\ 2,\ 9,\ 3$} 
        \textcolor{red}{\ealkey{Mean} $=5,\ \text{\ealkey{median} }=4$}
  \item \ealpara{\ealkeytr{Średnia}, \ealkeytr{mediana} i \ealkeytr{dominanta} z $5,\ 3,\ 7,\ 3,\ 5,\ 5,\ 9$}{\ealkey{Mean}, \ealkey{median} and \ealkey{mode} of $5,\ 3,\ 7,\ 3,\ 5,\ 5,\ 9$} 
        \textcolor{red}{\ealkey{Mean} $=\frac{37}{7}$, \ealkey{median} $=5$, \ealkey{mode} $=5$}
  \item \ealpara{Tylko \ealkeytr{dominanta} to $4$}{Only \ealkey{mode} is $4$} 
        \textcolor{red}{$x=4$}
\end{enumerate}
\end{frame}


\begin{frame}{Starter}
\vspace{-0.8em}

\footnotesize
\begin{enumerate}\itemsep2pt
\begin{multicols}{2}
  \item $3+4 = \ashowq{7}$
  \item $4 \times 8 = \ashowq{32}$
  \item $-3 + 8 + (-5) = \ashowq{0}$
  \item $\frac{2}{3} + \frac{1}{6} = \ashowq{\frac{5}{6}}$
  \item $-12 \div 4 = \ashowq{-3}$
  \item $\frac{3}{4} \times 8 = \ashowq{6}$
  \end{multicols}
  \item \ealpara{Klasa głosuje na swój ulubiony kolor. Czy możesz znaleźć \ealkeytr{średnią}? Której \ealkeytr{średniej} powinieneś użyć zamiast tego i dlaczego?}{A class votes for their favourite colour. Can you find the \ealkey{mean}? Which \ealkey{average} should you use instead, and why?} 
        \ealpara{Nie; użyj \ealkeytr{dominanty}, ponieważ \ealkeytr{dane} są kategoriami.}{No; use the \ealkey{mode}, because the \ealkey{data} are categories.}
        \ashow{No; use the \ealkey{mode}, because the \ealkey{data} are categories.}
  \item \ealpara{Znajdź \ealkeytr{średnią}, \ealkeytr{medianę} i \ealkeytr{dominantę} z:}{Find the \ealkey{mean}, \ealkey{median} and \ealkey{mode} of:} $-3,\ 1,\ -3,\ 4,\ -1,\ 2,\ -3$ 
        \ealpara{\ealkeytr{Średnia} $=-\frac{3}{7}$, \ealkeytr{mediana} $=-1$, \ealkeytr{dominanta} $=-3$.}{\ealkey{Mean} $=-\frac{3}{7}$, \ealkey{median} $=-1$, \ealkey{mode} $=-3$.}
        \ashow{\ealkey{Mean} $=-\frac{3}{7}$, \ealkey{median} $=-1$, \ealkey{mode} $=-3$.}
  \item \ealpara{Znajdź \ealkeytr{medianę} z:}{Find the \ealkey{median} of:} $\frac{1}{4},\ \frac{3}{4},\ \frac{1}{2},\ \frac{1}{8},\ \frac{5}{8}$ \ashow{$\frac{1}{2}$}
  \item \ealpara{Wieki w klubie młodzieżowym to: $11,\ 12,\ 12,\ 13,\ 13,\ 13,\ 14,\ 42$ Której \ealkeytr{średniej} nie użyłbyś w tym przykładzie i dlaczego?}{The ages in a youth club are: $11,\ 12,\ 12,\ 13,\ 13,\ 13,\ 14,\ 42$ Which \ealkey{average} would you not use for this example, and why?} 
        \ealpara{\ealkeytr{Średniej}, ponieważ $42$ jest \ealkeytr{wartością odstającą}.}{\ealkey{Mean}, because $42$ is an \ealkey{outlier}.}
        \ashow{\ealkey{Mean}, because $42$ is an \ealkey{outlier}.}
  \item \ealpara{\ealkeytr{Średnia} z $\{x,\ 3,\ x,\ 7,\ 5\}$ wynosi $6$. Znajdź $x$ i podaj \ealkeytr{dominantę}.}{The \ealkey{mean} of $\{x,\ 3,\ x,\ 7,\ 5\}$ is $6$. Find $x$ and state the \ealkey{mode}.} 
        \ealpara{$x=6$, \ealkeytr{dominanta} $=6$}{$x=6$, \ealkey{mode} $=6$}
        \ashow{$x=6$, \ealkey{mode} $=6$}
  \item \ealpara{Pięć liczb ma \ealkeytr{średnią} $8$, \ealkeytr{medianę} $7$ i \ealkeytr{jedyną} \ealkeytr{dominantę} $5$. Dwie liczby to $5$ i $5$. Znajdź możliwy \ealkeytr{zbiór}.}{Five numbers have \ealkey{mean} $8$, \ealkey{median} $7$ and \ealkey{unique} \ealkey{mode} of $5$. Two numbers are $5$ and $5$. Find a possible \ealkey{set}.} 
        \ashow{$\{5,5,6,7,17\}$}
\end{enumerate}

\end{frame}


\begin{frame}{Starter -- Solutions}
\vspace{-0.8em}

\footnotesize
\begin{enumerate}\itemsep2pt
\begin{multicols}{2}
  \item $3+4 = \mathord{?}$ \textcolor{red}{$7$}
  \item $4 \times 8 = \mathord{?}$ \textcolor{red}{$32$}
  \item $-3 + 8 + (-5) = \mathord{?}$ \textcolor{red}{$0$}
  \item $\frac{2}{3} + \frac{1}{6} = \mathord{?}$ \textcolor{red}{$\frac{5}{6}$}
  \item $-12 \div 4 = \mathord{?}$ \textcolor{red}{$-3$}
  \item $\frac{3}{4} \times 8 = \mathord{?}$ \textcolor{red}{$6$}
\end{multicols}
  \item \ealpara{Czy możesz znaleźć \ealkeytr{średnią}?}{Can you find the \ealkey{mean}?} 
        \ealpara{Nie — \ealkeytr{dane} są kategoryczne}{No — \ealkey{data} is categorical} 
        \textcolor{red}{No — \ealkey{data} is categorical}
  \item \ealpara{\ealkeytr{Średnia}, \ealkeytr{mediana} i \ealkeytr{dominanta}}{\ealkey{Mean}, \ealkey{median} and \ealkey{mode}} 
        \textcolor{red}{\ealkey{Mean} $=-\frac{3}{7}$, \ealkey{median} $=-1$, \ealkey{mode} $=-3$}
  \item \ealpara{\ealkeytr{Mediana}}{\ealkey{Median}} 
        \textcolor{red}{$=\frac{1}{2}$}
  \item \ealpara{Której \ealkeytr{średniej} nie użyć i dlaczego?}{Which \ealkey{average} not to use and why?} 
        \ealpara{\ealkeytr{Średniej}, ponieważ 42 jest \ealkeytr{wartością odstającą}}{\ealkey{Mean}, because 42 is an \ealkey{outlier}} 
        \textcolor{red}{\ealkey{Mean}, because 42 is an \ealkey{outlier}}
  \item \ealpara{$x=?$ i \ealkeytr{dominanta}}{$x=?$ and \ealkey{mode}} 
        \textcolor{red}{$x=6$, \ealkey{mode} $=6$}
  \item \ealpara{Możliwy \ealkeytr{zbiór}}{Possible \ealkey{set}} 
        \textcolor{red}{$\{5,5,6,7,17\}$}
\end{enumerate}
\end{frame}


\begin{frame}{Starter}
\vspace{-0.8em}

\footnotesize
\begin{enumerate}\itemsep2pt
\begin{multicols}{2}
\item $3+8 = \ashowq{11}$
\item $120 \div 10 = \ashowq{12}$
  \item $-6 + (-2) + 5 = \ashowq{-3}$
  \item \ealpara{Ustaw od najmniejszej do największej:}{Order from smallest to largest:} $\frac{1}{3},\ \frac{5}{6},\ \frac{1}{2},\ \frac{2}{3}$ \ashow{$\frac{1}{3},\ \frac{1}{2},\ \frac{2}{3},\ \frac{5}{6}$}
  \item $3\frac{1}{2} + 2\frac{3}{4} = \ashowq{6\frac{1}{4}}$
  \item $-20 \div 4 = \ashowq{-5}$
  \end{multicols}

  \item \ealpara{Znajdź \ealkeytr{średnią}, \ealkeytr{medianę} i \ealkeytr{dominantę} z:}{Find the \ealkey{mean}, \ealkey{median} and \ealkey{mode} of:} $-6,\ -2,\ 4,\ -3,\ 7,\ 0,\ -2$ 
        \ealpara{\ealkeytr{Średnia} $=-\frac{2}{7}$, \ealkeytr{mediana} $=-2$, \ealkeytr{dominanta} $=-2$.}{\ealkey{Mean} $=-\frac{2}{7}$, \ealkey{median} $=-2$, \ealkey{mode} $=-2$.}
        \ashow{\ealkey{Mean} $=-\frac{2}{7}$, \ealkey{median} $=-2$, \ealkey{mode} $=-2$.}

  \item \ealpara{Prawda czy fałsz? „Dodanie 10 do każdej \ealkeytr{wartości} zwiększa \ealkeytr{średnią} o 10, ale nie zmienia \ealkeytr{mediany}.”}{True or False? ``Adding 10 to every \ealkey{value} increases the \ealkey{mean} by 10 but leaves the \ealkey{median} unchanged.''} 
        \ealpara{Prawda}{True}
        \ashow{True}

  \item \ealpara{Siedem \ealkeytr{wartości} $n-9,\ n-6,\ n-3,\ n,\ n+3,\ n+6,\ n+9$ ma \ealkeytr{średnią} $4$. Podaj \ealkeytr{medianę} \textbf{bez obliczeń}.}{The seven \ealkey{values} $n-9,\ n-6,\ n-3,\ n,\ n+3,\ n+6,\ n+9$ have a \ealkey{mean} of $4$. State the \ealkey{median} \textbf{without calculation}.} 
        \ealpara{\ealkeytr{Mediana} $=4$}{\ealkey{Median} $=4$}
        \ashow{\ealkey{Median} $=4$}
  \item \ealpara{\ealkeytr{Zbiór} pięciu \ealkeytr{dodatnich} \ealkeytr{liczb całkowitych} ma \ealkeytr{średnią} = \ealkeytr{medianę} = \ealkeytr{dominantę} = $6$.}{A \ealkey{set} of five \ealkey{positive} \ealkey{integers} has \ealkey{mean} = \ealkey{median} = \ealkey{mode} = $6$.}
        \begin{itemize}\footnotesize
          \item \ealpara{Jaki jest najmniejszy możliwy \ealkeytr{rozstęp}?}{What is the smallest possible \ealkey{range}?} \ashow{$1$}
          \item \ealpara{Jaki jest największy możliwy \ealkeytr{rozstęp}?}{What is the largest possible \ealkey{range}?} \ashow{$12$}
        \end{itemize}
\end{enumerate}

\end{frame}


\begin{frame}{Starter -- Solutions}
\vspace{-0.8em}

\footnotesize
\begin{enumerate}\itemsep2pt
\begin{multicols}{2}
  \item $3+8 = \mathord{?}$ \textcolor{red}{$11$}
  \item $120 \div 10$ \textcolor{red}{$12$}
  \item $-6 + (-2) + 5$ \textcolor{red}{$-3$}
  \item \ealpara{Ustaw od najmniejszej do największej:}{Order from smallest to largest:} 
        \textcolor{red}{$\frac13,\ \frac12,\ \frac23,\ \frac56$}
  \item $3\frac12 + 2\frac34$ \textcolor{red}{$6\frac14$}
  \item $-20 \div 4$ \textcolor{red}{$-5$}
\end{multicols}
  \item \ealpara{\ealkeytr{Średnia}, \ealkeytr{mediana}, \ealkeytr{dominanta}}{\ealkey{Mean}, \ealkey{median}, \ealkey{mode}} 
        \textcolor{red}{\ealkey{Mean} $=-\frac{2}{7}$, \ealkey{median} $=-2$, \ealkey{mode} $=-2$}
  \item \ealpara{Prawda czy fałsz?}{True or False?} 
        \ealpara{Prawda}{True}
        \textcolor{red}{True}
  \item \ealpara{\ealkeytr{Mediana}}{\ealkey{Median}} 
        \textcolor{red}{$4$}
  \item \ealpara{\ealkeytr{Rozstęp}}{\ealkey{Range}} 
        \ealpara{Najmniejszy $=1$, największy $=12$}{Smallest $=1$, largest $=12$} 
        \textcolor{red}{Smallest $=1$, largest $=12$}
\end{enumerate}
\end{frame}

\end{document}